\part{Economic Activity and the Establishment of Futures}
\label{part:I}

Part I fixes what an economy maintains and introduces the objects the rest of the book relies on: states, histories, continuations, and the distinction between what is possible, accessible, permitted, and acceptable.

\begin{housing}Define the building, its tenants, and the repair crew, and list the states and histories that matter: rent arrears, deferred work, and notices served.\end{housing}

\chapter{What an Economy Maintains}\label{ch:1}
Introduces an economy as an organization of provisioning, transformation, and maintenance through time, using a household, a construction firm, a hospital, and a watershed to show the range of relevant processes. Separates current output from the reproduction of the conditions that make output possible, and shows how one transaction can generate income while exhausting productive capacity elsewhere.
\scaffold

\chapter{States, Histories, and Continuations}\label{ch:2}
Introduces economic states, histories, actions, and transition relations. States when a present-state description is sufficient and when accumulated obligations, deterioration, learning, or institutional commitments require history to remain explicit. The basic object is a set of possible trajectories, not a single forecast. The first theorem characterizes sufficient state augmentation: a compressed history is adequate when histories with the same summary have equivalent consequences for the future questions at issue.
\[
x_{t+1}\in T(x_t,h_t,a_t),\qquad h_{t+1}=U(h_t,x_t,a_t,x_{t+1}).
\]
\chspec{\Pref{P01}.}
\scaffold

\chapter{Feasibility, Accessibility, and Admissibility}\label{ch:3}
Separates physically possible actions, practically accessible actions, institutionally permitted actions, and actions consistent with stated obligations. Examines why formal eligibility does not establish effective access and why an available choice need not be an acceptable continuation. Examples are housing, employment, medical treatment, and access to credit.
\chspec{\Pref{P08}.}
\scaffold

\chapter{Establishing the Domain of Choice}\label{ch:4}
Examines the work required to make alternatives available before selection can occur. Search, education, infrastructure, legal recognition, invention, and interface design become explicit economic activities. A two-stage model distinguishes expenditures that establish an opportunity domain from expenditures that select within it.
\chspec{\Pref{P02}.}
\scaffold

\chapter{Welfare Without Counting Options}\label{ch:5}
Addresses the normative foundations directly. More possible trajectories do not automatically make a better life, and constraints can protect some people while excluding others. Develops plural assessments involving substantive achievement, security, burden, access, and corrective freedom, and makes disagreements about those assessments explicit.
\scaffold
